\subsection{Geometric Properties of affine Springer fibers}


\subsubsection{Non-reducedness} The ind-scheme $\tX_{\gamma}$ is {\em never} reduced if $G$ is nontrivial and $\gamma$ is regular semisimple in $\frg(F)$. For example, in the case considered in \S\ref{sss:reduction rs}, $\tX_{\gamma}$ is isomorphic to $\Gr_{T}=LT/L^{+}T$. We have seen in \S\ref{sss:centralizer} that for non-reduced rings $R$, elements in $LT(R)=R\lr{t}^{\times}$ can have nilpotent leading coefficients. Therefore $\Gr_{T}(R)$ is not just $\xcoch(T)$, which is $\Gr^{\red}_{T}$. This shows that $\Gr_{T}$ is non-reduced, hence $\tX_{\gamma}$ is non-reduced. 


The next theorem is the fundamental finiteness statement about $\sX_{\gamma}$.
 
 
 
 
\begin{theorem}[Essentially Kazhdan and Lusztig {\cite[Prop 2.1]{KL}}]\label{th:ft quot} Let $\gamma\in \frg(F)$ be a regular semisimple element. Then
\begin{enumerate}
\item There exists a closed subscheme $Z\subset \sX_{\gamma}$ which is projective over $k$, such that $\sX_{\gamma}=\cup_{\ell\in\L_{\gamma}}\ell\cdot Z$.
\item The ind-scheme $\sX_{\gamma}$ is a scheme locally of finite type over $k$. 
\item The action of $\L_{\gamma}$ on $\sX_{\gamma}$ is free, and the quotient $\L_{\gamma}\bs\sX_{\gamma}$ (as an fppf sheaf on $k$-algebras) is representable by a proper algebraic space over $k$.
\end{enumerate}
\end{theorem}

We sketch a proof of this theorem below in three steps.

\subsubsection{First reduction} We show that part (1) of the theorem implies (2) and (3). Let $Z$ be a projective subscheme as in (1). To show (2), we would like to show that any $x\in \sX_{\gamma}(k)$ has an open neighborhood which is a scheme of finite type. By the $\L_{\gamma}$-action we may assume $x\in Z(k)$. Since $Z$ is of finite type, the set $\Sigma:=\{\ell\in \L_{\gamma}|Z\cap\ell\cdot Z\neq\varnothing\}$ is finite. Let $U=\sX_{\gamma}-\cup_{\ell\notin\Sigma}\ell\cdot Z$, then $U$ is an open neighborhood of $Z$, hence an open  neighborhood of $x$. Moreover, $U$ is contained in the finite union $\cup_{\ell\in\Sigma}\ell\cdot Z$, hence contained in some Schubert variety $\Gr_{G,\leq\l}$. Hence $U$ is an open subset of the projective scheme $\sX_{\gamma}\cap\Gr_{G,\leq\l}$, therefore $U$ is itself a scheme of finite type. To show (3), note that the fppf sheaf quotient $\L_{\gamma}\bs\sX_{\gamma}$ is a separated algebraic space because it is the quotient of $\sX_{\gamma}$ by the \'etale equivalence relation $\L_{\gamma}\times\sX_{\gamma}\subset \sX_{\gamma}\times\sX_{\gamma}$ (given by the action and projection maps). By (1), there is a surjection $Z\to \L_{\gamma}\bs\sX_{\gamma}$ from a projective scheme $Z$, which implies that $\L_{\gamma}\bs\sX_{\gamma}$ is proper.

\subsubsection{Proof of (1) when $\gamma$ lies in a split torus}\label{sss:split case} We first consider the case where $\gamma$ lies in a split torus. By $G(F)$-conjugation, we may assume $\gamma\in\frt(F)$. In this case $\L_{\gamma}\cong\xcoch(T)$. Fix a Borel subgroup $B$ containing $T$ and let $N$ be the unipotent radical of $B$. The Iwasawa decomposition of $LG$ gives
\begin{equation*}
\Gr_{G}=LN\cdot \L_{\gamma}L^{+}G/L^{+}G=\sqcup_{\l\in\xcoch(T)}LN\cdot \l(t)L^{+}G/L^{+}G.
\end{equation*}
Let $X:=\sX_{\gamma}\cap (LN\cdot L^{+}G/L^{+}G)\subset \sX_{\gamma}$. It is enough to show that $X$ lies in some affine Schubert variety $\Gr_{\leq \l}$, for then its closure $Z:=\ov{X}$ in $\Gr_{\leq\l}$ satisfies the condition in (1). For later use, we note that the translations $\ell\cdot X$ for $\ell\in \L_{\gamma}$ are disjoint and cover $\sX_{\gamma}$. 

Fixing an ordering of the positive roots of $T$ with respect to $B$, we may write an element  $u\in N(F)$ uniquely as
\begin{equation}\label{root gp}
u=\prod_{\alpha>0}x_{\alpha}(c_{\alpha})
\end{equation}
where $c_{\alpha}\in F$ and $x_{\alpha}: \Ga\to N$ is the root group corresponding to $\alpha$. To show that $X$ is contained in an affine Schubert variety, it suffices to give a lower bound for the valuations of $c_{\alpha}$ appearing in \eqref{root gp} for any $u\in N(F)$ such that $[u]\in X(k)$. We may expand $\Ad(u^{-1})\gamma\in\frb(F)$ in terms of the root decomposition
\begin{equation*}
\Ad(u^{-1})\gamma=\gamma+\sum_{\alpha>0}P_{\alpha}(\gamma; c)e_{\alpha}
\end{equation*} 
where $e_{\alpha}\in\frg_{\alpha}$ is a fixed basis for each root space and $P_{\alpha}(\gamma;c)$ is an $F$-valued polynomial function in $\{c_{\alpha}\}_{\alpha>0}$ and linear in $\gamma$. Induction on the height of $\alpha$ shows that $P_{\alpha}(\gamma;c)$ takes the following form
\begin{equation}\label{conj u gamma}
P_{\alpha}(\gamma;c)=\jiao{\alpha,\gamma}c_{\alpha}+\sum_{\beta<\alpha}\jiao{\beta,\gamma}P^{\beta}_{\alpha}(c)
\end{equation}
where $P^{\beta}_{\alpha}(c)$ is a polynomial involving only $\{c_{\alpha'}\}_{\alpha'<\alpha}$, and homogeneous of degree $\alpha$ (we define $\deg c_{\alpha'}:=\alpha'\in\xch(T)$). 

Let $n=\max_{\alpha>0}\{\val_{F}\jiao{\alpha,\gamma}\}$. This is finite because $\gamma$  is regular semisimple. If $[u]\in X(k)$, i.e., $\Ad(u^{-1})\gamma\in\frg(\calO_{F})$, then induction on the height of $\alpha$ shows that
\begin{equation*}
\val_{F}(c_{\alpha})\geq -(2\textup{ht}(\alpha)-1)n,
\end{equation*} 
which gives the desired lower bound and shows that $X$ lies in an affine Schubert variety.

\subsubsection{Proof of (1) in the general case} In the general case, we give a simplified argument compared to the original one in \cite{KL}, following the same idea. We make a base change to $F'=k\lr{t^{1/m}}$ over which $\gamma$ can be conjugated into a split torus. Let $\Gr'_{G}$ be the affine Grassmannian of $G$ defined using the field $F'$ in place of $F$ (so that $\Gr'_{G}(k)=G(F')/G(\calO'_{F})$), and let $\sX'_{\gamma}\subset\Gr'_{G}$ be the corresponding affine Springer fiber. Then both $\Gr'_{G}$ and $\sX'_{\gamma}$ carry an action of $\Gamma:=\Gal(F'/F)\cong\mu_{m}$ induced from its action on $F'$, and we have a closed embedding $\sX_{\gamma}\incl (\sX'_{\gamma})^{\Gamma}$. 

Let $\L'_{\gamma}=\xcoch(G_{\gamma}\otimes F')\incl G_{\gamma}(F')$ be the lattice constructed using the field $F'$. There is an action of $\Gamma$ on $\xcoch(G_{\gamma}\otimes F')$ with fixed points $\xcoch(G_{\gamma})\cong \L_{\gamma}$. This action induces an action of $\Gamma$ on $\L'_{\gamma}$, but it does not respect the embedding $\xcoch(G_{\gamma}\otimes F')\incl G_{\gamma}(F')$. The fixed points $(\L'_{\gamma})^{\Gamma}=\xcoch(G_{\gamma}\otimes F')^{\Gamma}$ may not lie in $G_{\gamma}(F)$, however it always contains $\L_{\gamma}$ with finite index. 

From the proof in the split case in \S\ref{sss:split case} we have a finite type locally closed subscheme $X'\subset\sX'_{\gamma}$ coming from the Iwasawa decomposition, such that $\sX'_{\gamma}$ can be decomposed as the {\em disjoint union} $\sqcup_{\ell\in\L'_{\gamma}}\ell\cdot X'$  (not as schemes but as constructible sets). We identify $\ell\in \L'_{\gamma}$ with an element in $\xcoch(G_{\gamma}\otimes F')$ then $\sigma(\ell)$ makes sense for $\sigma\in\Gamma$. One checks that $\sigma(\ell\cdot X')=\sigma(\ell)\cdot X'$ for $\ell\in\L'_{\gamma}$ even though $\sigma$ does not respect the embedding $\L'_{\gamma}\to G_{\gamma}(F')$. Therefore $(\sX'_{\gamma})^{\Gamma}\subset (\L'_{\gamma})^{\Gamma}\cdot X'$. Choosing representatives $C$ for the finite coset space $(\L'_{\gamma})^{\Gamma}/\L_{\gamma}$, we see that $(\sX'_{\gamma})^{\Gamma}$ is contained in $\L_{\gamma}\cdot (C\cdot X')$. Since $C\cdot X'$ is of finite type, $Z'=\ov{C\cdot X'}\cap (\sX'_{\gamma})^{\Gamma}$ is a projective subscheme of $\sX'_{\gamma}$ whose $\L_{\gamma}$-translations cover $(\sX'_{\gamma})^{\Gamma}$. Finally the projective subscheme $Z=Z'\cap \sX_{\gamma}$ of $\sX_{\gamma}$ satisfies the requirement of (1).  This finishes the proof of Theorem \ref{th:ft quot}. \qed


\subsubsection{Reduction to Levi} The proof of Theorem \ref{th:ft quot} in the split case in \S\ref{sss:split case} gives more information. In the Iwasawa decomposition, let $S_{\l}\subset \Gr_{G}$ be the $LN$-orbit of $\l(t)$, for $\l\in\xcoch(T)$. This is called a {\em semi-infinite orbit}, because it has infinite dimension and also has infinite codimension in $\Gr_{G}$.  Let $C_{\l}:=\sX_{\gamma}\cap S_{\l}$, then $C_{\l}=\l(t)\cdot X$ in the notation of \S\ref{sss:split case}. The formula \eqref{conj u gamma} implies that $C_{\l}\neq\varnothing$ (or equivalently $X\neq\varnothing$, or $\sX_{\gamma}\neq\varnothing$) if and only if $\jiao{\alpha,\gamma}\in\calO_{F}$ for all roots $\alpha$, and if so, $C_{\l}$ is isomorphic to an almost affine space (namely an iterated $\AA^{1}$-bundle) of dimension
\begin{equation*}
\dim C_{\l}=\sum_{\alpha>0}\val_{F}(\jiao{\alpha,\gamma})=\frac{1}{2}\val_{F}\Delta(\gamma)
\end{equation*}
Here $\Delta(\gamma)$ is the determinant of the adjoint action of $\gamma$ on $\frg(F)/\frt(F)$. Therefore $\sX_{\gamma}$ can be decomposed into almost affine spaces of the same dimension indexed by $\xcoch(T)$. However, this decomposition is not a stratification: the closure $\ov{C}_{\l}$ of $C_{\l}$ will intersect other $C_{\l'}$ but certainly not a union of such $C_{\l'}$'s. 

The decomposition $\sX_{\gamma}=\sqcup C_{\l}$ in the split case has a generalization. Suppose $P$ is a parabolic subgroup of $G$ with unipotent radical $N_{P}$ and a Levi subgroup $M_{P}$.  Let $\mathfrak{m}_{P}=\Lie M_{P}$ and suppose $\gamma\in\mathfrak{m}_{P}(F)$ is regular semisimple as an element in  $\frg(F)$.  Using the generalized Iwasawa decomposition $G(F)=N_{P}(F)M_{P}(F)G(\calO_{F})$, there is a well-defined map $\Gr_{G}(k)\to \Gr_{M_{P}}(k)$ sending $nm G(\calO_{F})$ to $m M_{P}(\calO_{F})$, for $n\in N_{P}(F)$ and $m\in M_{P}(F)$. However this map does not give a map of ind-schemes. Nevertheless the fibers of this map have natural structure of infinite dimensional affine spaces. Restricting this map to $\sX_{\gamma}$ we get
\begin{equation*}
\tau: \sX_{\gamma}(k)\to \sX^{M_{P}}_{\gamma}(k)
\end{equation*}
where $\sX^{M_{P}}_{\gamma}\subset \Gr_{M_{P}}$ is the affine Springer fiber for $\gamma$ and the group $M_{P}$. The fibers of $\tau$, if non-empty, are almost affine spaces of dimension $\frac{1}{2}\val_{F}\Delta^{G}_{M_{P}}(\gamma)$, where
\begin{equation}\label{Delta G L}
\Delta^{G}_{M_{P}}(\gamma):=\det(\ad(\gamma)|\frg(F)/\mathfrak{m}_{P}(F)).
\end{equation}
Assume $T\subset M_{P}$, then the connected components of $\Gr_{M_{P}}$ are indexed by $\xcoch(T)/R^{\vee}_{M_{P}}$ where $R^{\vee}_{M_{P}}$ is the coroot lattice of $M_{P}$. If we decompose $\sX^{M_{P}}_{\gamma}$ into connected components $\sX^{M_{P}}_{\gamma}(\l)$ for $\xcoch(T)/R^{\vee}_{M_{P}}$ and taking their preimages under $\tau$, we get a decomposition $\sX_{\gamma}=\sqcup\sX_{\gamma,\l}$ into locally closed sub-ind-schemes indexed also by $\l\in\xcoch(T)/R^{\vee}_{M_{P}}$. One can show that $\sX_{\gamma,\l}$, if non-empty, is an almost affine space bundle over $\sX^{M_{P}}_{\gamma}(\l)$ with fiber dimension $\frac{1}{2}\val_{F}\Delta^{G}_{M_{P}}(\gamma)$.



\subsubsection{Connectivity and equidimensionality}\label{sss:equidim} When $G$ is simply-connected,  $\Fl$ is connected, and in this case the affine Springer fiber $\sY_{\gamma}$ is also connected. See \cite[\S4, Lemma 2]{KL}. As a consequence, when $G$ is simply-connected,  $\sX_{\bP,\gamma}$ is connected for all parahoric $\bP$ because the natural projection $\sY_{\gamma}\to \sX_{\bP,\gamma}$ is surjective. 

In \cite{KL}, it is also shown that $\sY_{\gamma}$ is equidimensional. The argument there is similar to Spaltenstein's the proof of the connectivity and equidimensionality for Springer fibers in \cite{Spa77}.


\subsubsection{The dimension formula}
By Theorem \ref{th:ft quot}, the dimension of $\sX_{\gamma}$ is well-defined, and is the dimension of $L_{\gamma}\bs\sX_{\gamma}$ as an algebraic space. To state a formula for $\dim\sX_{\gamma}$, we need some more notation.

Consider the adjoint action $\ad(\gamma):\frg(F)\to \frg(F)$. The kernel of this map is $\frg_{\gamma}(F)$, and the induced endomorphism $\ov{\ad(\gamma)}$ on $\frg(F)/\frg_{\gamma}(F)$ is invertible. Let $\Delta(\gamma)\in F^{\times}$ be the determinant of $\ov{\ad(\gamma)}$. This is consistent with our earlier definition of $\Delta(\gamma)$ in the case $\gamma$ lies in a split torus $\frt(F)$.

On the other hand, recall $\xcoch(G_{\gamma})$ is the group of $F$-rational cocharacters of $G_{\gamma}$, which is also the rank of the maximal $F$-split subtorus of $G_{\gamma}$. Let
\begin{equation*}
c(\gamma)=r-\rk_{\ZZ}\xcoch(G_{\gamma}).
\end{equation*}
Then $c(\gamma)$ is also the rank of the maximal $F$-anisotropic subtorus $G_{\gamma}$.

\begin{theorem}[Bezrukavnikov {\cite{Bez}}, conjectured by Kazhdan-Lusztig \cite{KL}]\label{th:aff dim} Let $\gamma\in \frg(F)$ be a regular semisimple element. Then we have 
\begin{equation}\label{aff Spr dim}
\dim\sX_{\gamma}=\frac{1}{2}(\val_{F}\Delta(\gamma)-c(\gamma)).
\end{equation}
\end{theorem}



\subsubsection{Sketch of proof}\label{dim formula discussion} A key role in the proof is played by the notion of {\em regular points} of $\sX_{\gamma}$. We have an evaluation map $\ev: \sX_{\gamma}\to [\frg/G]$ sending $[g]\in\sX_{\gamma}\subset \Gr_{G}$ to the reduction of $\Ad(g^{-1})\gamma$ modulo $t$, which is well-defined up to the adjoint action by $G$.  We say $[g]\in \sX_{\gamma}$ is a {\em regular point} if $\ev([g])$ lies in the open substack $[\frg^{\reg}/G]$ of $[\frg/G]$. Let $\sX^{\reg}_{\gamma}\subset\sX_{\gamma}$ be the open sub-ind-scheme of $\sX_{\gamma}$ consisting of regular points. It can be shown that $\sX^{\reg}_{\gamma}$ is non-empty. Denote the preimage of $\sX^{\reg}_{\gamma}$ in $\sY_{\gamma}$ by $\sY^{\reg}_{\gamma}$, then the projection map $\sY^{\reg}_{\gamma}\to\sX^{\reg}_{\gamma}$ is an isomorphism. Since $\sY^{\reg}_{\gamma}$ is equidimensional as mentioned in \S\ref{sss:equidim}, we see that $\dim\sX^{\reg}_{\gamma}=\dim\sY^{\reg}_{\gamma}=\dim\sY_{\gamma}$. Of course we have $\dim\sX_{\gamma}\leq\dim \sY_{\gamma}$, therefore we must have $\dim \sX^{\reg}_{\gamma}=\dim \sX_{\gamma}=\dim \sY_{\gamma}$. It remains to calculate the dimension of $\sX^{\reg}_{\gamma}$.

Recall we have defined the local Picard group $P_{a}$ in \S\ref{sss:sym ASF}. The action of $P_{a}$ on $\sX_{\gamma}$ preserves $\sX^{\reg}_{\gamma}$, and in fact $\sX^{\reg}_{\gamma}$ is a torsor under $P_{a}$.  Therefore it suffices to compute the dimension of $P_{a}=LG_{\gamma}/L^{+}J_{a}$. 

Consider the projection $\chi_{\frt}: \frt\to\frc=\frt\sslash W$. Pullback along $a=\chi(\gamma):\Spec\calO_{F}\to\frc$ we get a finite morphism $\chi_{a}: \chi^{-1}_{\frt}(a)=\Spec A\to \Spec\calO_{F}$, for a finite flat $\calO_{F}$-algebra $A$. There is a close relationship between the group scheme $J_{a}$ and $J'_{a}=(\Res^{A}_{\calO_{F}}(T\otimes_{k}A))^{W}$, where the Weyl group $W$ acts diagonally on both $T$ and $A$.  In fact $J_{a}$ and $J'_{a}$ are equal up to connected components in their special fibers. In particular, $\dim P_{a}=\dim P'_{a}$, where $P'_{a}=LG_{\gamma}/L^{+}J'_{a}$. It is not hard to see that
\begin{equation*}
\dim P'_{a}=\dim (\frt\otimes \tilA)^{W}-\dim (\frt\otimes A)^{W}=\dim (\frt\otimes (\tilA/A))^{W}
\end{equation*}
where $\tilA$ is the normalization of $A$. From this one deduces the dimension formula \eqref{aff Spr dim}. \qed

\begin{remark} Using the relation between affine Springer fibers and Hitchin fibers, Ng\^o \cite[Cor 4.16.2]{NgoFL} showed that $\sX^{\reg}_{\gamma}$ is in fact dense in $\sX_{\gamma}$. In particular, each irreducible component of $\sX_{\gamma}$ is a rational variety. However, this rationality property is false for affine Springer fibers in more general affine partial flag varieties, as we saw in Bernstein-Kazhdan's example in \S\ref{sss:BK example}. 
\end{remark}


\subsubsection{Purity} It was conjectured by Goresky, Kottwitz and MacPherson \cite{GKM} that the cohomology of affine Springer fibers should be pure (in the sense of Frobenius weights if $k=\overline{\FF}_{p}$, or in the sense of Hodge structures if $k=\CC$). The purity of affine Springer fibers would allow the authors of \cite{GKM} to prove the Fundamental Lemma for unramified elements using localization techniques in equivariant cohomology. This purity conjecture is still open in general. In \cite{GKMPurity}, a class of affine Springer fibers called {\em equivalued} were shown to be pure.


\subsubsection{Invariance under perturbation}  Suppose $a,a'\in\frc(\calO_{F})\cap\frc^{\rs}(F)$. We say $a\equiv a'\mod t^{N}$ if $a$ and $a'$ have the same image under the map $\frc(\calO_{F})\to\frc(\calO_{F}/t^{N})$. In \cite[Prop 3.5.1]{NgoFL}, it is shown that for fixed $a\in\frc(\calO_{F})\cap\frc^{\rs}(F)$, there exists some $N$ (depending on $a$) such that whenever $a'\equiv a\mod t^{N}$, we have isomorphisms
\begin{equation*}
\tX_{a}\cong\tX_{a'} \textup{ and } P_{a}\cong P_{a'}
\end{equation*}
in a way compatible with the actions. Therefore, we may say that $\sX_{a}$ varies locally constantly with $a$ under the $t$-adic topology on $\frc(\calO_{F})$. 

For example,  consider the case $G=\GL_{n}$ and let $a\in\frc(\calO_{F})$ correspond to a characteristic polynomial $P(x)=x^{n}+a_{1}x^{n-1}+\cdots+a_{n}$ whose roots are in $\calO_{F}$ and are distinct modulo $t$. Then for any $a'\equiv a\mod t$, the characteristic polynomial of $a'$ also has distinct roots modulo $t$. In this case, $\tX_{a}$ and $\tX_{a'}$ are both torsors under $P_{a}\cong P_{a'}\cong \Gr_{T}$.



